\section*{Capítulo \ref{ch:b1:complexation} --- Complexação}

\begin{solution}{exo:b1:complexation:1}
$[\ce{Cu(NH3)4}]^{2+}$: \ce{Cu^2+}, quatro \ce{NH3} ligadas pelo N.
$[\ce{Fe(CN)6}]^{4-}$: \ce{Fe^2+}, seis \ce{CN-} ligados pelo C;
$[\ce{Fe(CN)6}]^{3-}$: o mesmo com \ce{Fe^3+}. $[\ce{Ag(NH3)2}]^{+}$:
\ce{Ag+}, duas \ce{NH3} pelo N. $[\ce{FeSCN}]^{2+}$: \ce{Fe^3+}, um
\ce{SCN-}, ligado pelo N apesar da maneira como a fórmula é escrita.
$[\ce{Zn(OH)4}]^{2-}$: \ce{Zn^2+}, quatro \ce{OH-} pelo O.
$[\ce{CaY}]^{2-}$: \ce{Ca^2+}, um íon edta por dois N e quatro O.
\end{solution}

\begin{solution}{exo:b1:complexation:2}
$\log K_i = \log\beta_i - \log\beta_{i-1}$: 4.10, 3.41, 2.82, 2.03, e
$\mathrm{p}K_{d,i} = \log K_i$ assume os mesmos valores. O último se refere a
$[\ce{Cu(NH3)4}]^{2+} \rightleftharpoons [\ce{Cu(NH3)3}]^{2+} + \ce{NH3}$.
\end{solution}

\begin{solution}{exo:b1:complexation:3}
Domínios: $[\ce{Cu(NH3)4}]^{2+}$ abaixo de pNH3 2.03, depois $n = 3$ até 2.82,
$n = 2$ até 3.41, $n = 1$ até 4.10, \ce{Cu^2+} acima. pNH3 = 0:
$[\ce{Cu(NH3)4}]^{2+}$; 2.5: $[\ce{Cu(NH3)3}]^{2+}$; 3.5:
$[\ce{CuNH3}]^{2+}$; 5.0: \ce{Cu^2+}.
\end{solution}

\begin{solution}{exo:b1:complexation:4}
$[\ce{FeSCN^2+}]/[\ce{Fe^3+}] = \beta_1[\ce{SCN-}] = 10^{2.96} \times 0.10 =
91$: a fração complexada é $91/92 = \qty{98.9}{\percent}$.
\end{solution}

\begin{solution}{exo:b1:complexation:5}
Cada $[\mathrm{ML}_k] = \beta_k[\mathrm{M}][\mathrm{L}]^k$; dividir pela
soma em $k$ (o metal total) elimina $[\mathrm{M}]$. Com três espécies,
$f_i = 1/(1 + [\mathrm{ML}_{i-1}]/[\mathrm{ML}_i] +
[\mathrm{ML}_{i+1}]/[\mathrm{ML}_i]) = 1/(1 + 1/(K_i x) + K_{i+1}x)$ com
$x = [\mathrm{L}]$. O denominador é mínimo quando sua derivada
$-1/(K_i x^2) + K_{i+1}$ se anula, $x^2 = 1/(K_iK_{i+1})$, onde as duas
vizinhas são iguais, e $\mathrm{pL} = \frac12(\log K_i + \log K_{i+1})$.
Para $[\ce{Cu(NH3)2}]^{2+}$: $\frac12(3.41 + 2.82) = 3.12$.
\end{solution}

\begin{solution}{exo:b1:complexation:6}
$[\ce{FeSCN}]^{2+} + \ce{F-} \rightleftharpoons [\ce{FeF}]^{2+} + \ce{SCN-}$,
$K = 10^{6.85 - 2.96} = 10^{3.89}$. Razão $= K[\ce{F-}]/[\ce{SCN-}] =
10^{3.89} \times 0.010/0.10 = 780$: o \omterm{def:b1:complexation:complex}{complexo} vermelho é quase todo
substituído e a cor desaparece.
\end{solution}

\begin{solution}{exo:b1:complexation:7}
$[\ce{MgY}]^{2-} + \ce{Ca^2+} \rightleftharpoons [\ce{CaY}]^{2-} + \ce{Mg^2+}$,
$K = 10^{12.69 - 10.90} = 10^{1.79} = 62$. Com quantidades iguais,
$x^2/(1 - x)^2 = 62$, $x/(1 - x) = 7.9$, $x = 0.89$: \qty{89}{\percent} do
magnésio é liberado.
\end{solution}

\begin{solution}{exo:b1:complexation:8}
\ce{Ag2O(s) + H2O + 4NH3 <=> 2[Ag(NH3)2]+ + 2OH-}, $K = \beta_2^2 \times
10^{-15.43} = 10^{14.44 - 15.43} = 10^{-0.99}$. A constante não é
pequena: assim que a amônia livre chega a alguns décimos de mol por
litro, o óxido marrom formado pelas primeiras gotas (a amônia agindo
como base) se dissolve no aminocomplexo de prata.
\end{solution}

\begin{solution}{exo:b1:complexation:9}
pH 10: $\alpha = 1 + 10^{1.24} + 10^{-1.96} = 18.4$, $\log\beta' = 12.69 -
1.26 = 11.43$. pH 7: $\alpha = 1 + 10^{4.24} + 10^{4.04} + 10^{0.19} +
\cdots = 10^{4.45}$, $\log\beta' = 8.24$. Em pH 7, a constante condicional
fica abaixo de $10^{10}$: a reação do cálcio com o edta não é completa o
bastante para uma titulação, porque a maior parte do edta está protonada.
\end{solution}

\begin{solution}{exo:b1:complexation:10}
1.~\ce{CuO(s) + H2O + 4NH3 <=> [Cu(NH3)4]^2+ + 2OH-}, $K = \beta_4 \times
10^{-20.64} = 10^{12.36 - 20.64} = 10^{-8.28}$.
2.~$\mathrm{pH} = 7 + \frac12(9.25 + \log 1.0) = 11.63$, $[\ce{OH-}] =
10^{-2.37}$, e $[\ce{Cu(NH3)4^2+}] = 10^{-8.28}/10^{-4.75} =
\qty{3.0e-4}{mol/L}$: muito pouco se dissolve.
3.~$\mathrm{pH} = 9.25$, $[\ce{OH-}] = 10^{-4.75}$, e a fórmula dá
$10^{-8.28 + 9.50} = 10^{1.22}$, muito mais do que a amônia pode
carregar: o óxido se dissolve até que a amônia, e não o equilíbrio, se
esgote. Os íons amônio consomem o hidróxido liberado pela dissolução; é
por isso que os sais de amônio ajudam a amônia a dissolver os compostos
de cobre.
\end{solution}

\begin{solution}{exo:b1:complexation:11}
1.~A reação é a segunda formação menos a primeira: $K = K_2/K_1 =
10^{3.90 - 3.32} = 10^{0.58} = 3.8 > 1$, de modo que $[\ce{AgNH3}]^{+}$
reage consigo mesmo quando é a espécie principal.
2.~Pelo \cref{exo:b1:complexation:5}, a maior fração ocorre em
$\mathrm{pNH_3} = \frac12(3.32 + 3.90) = 3.61$, onde $\beta_1 x = 10^{-0.29}
= 0.51$ e $\beta_2x^2 = 1.0$: $f = 0.51/(1 + 0.51 + 1.0) = 0.20$. Nunca
mais de \qty{20}{\percent} da prata está como $[\ce{AgNH3}]^{+}$.
\end{solution}

\begin{solution}{exo:b1:complexation:12}
1.~\ce{CaCO3(s) + Y^4- <=> [CaY]^2- + CO3^2-}, $K = \beta K_s = 10^{12.69 -
8.30} = 10^{4.39}$.
2.~$x^2/(0.10 - x) = 10^{4.39}$ dá $x = \qty{0.10}{mol/L}$ com precisão de
$4 \times 10^{-7}$: todo o edta é usado, \qty{10.0}{g} de carbonato de
cálcio por litro.
3.~Em meio ácido, o edta é protonado ($\mathrm{p}K_{a4} = 11.24$) e sua
constante condicional cai; o carbonato também é protonado, o que ajuda,
mas é o \omterm{def:b1:complexation:complex}{complexo} que retém o cálcio. Manter o pH alto mantém o edta como
\ce{Y^4-}.
\end{solution}

\begin{solution}{pb:b1:complexation:1}
\textbf{1.} \ce{Ag+ + NH3 <=> [AgNH3]+}, $\beta_1 = [\ce{AgNH3+}]/([\ce{Ag+}][\ce{NH3}])$;
\ce{Ag+ + 2NH3 <=> [Ag(NH3)2]+}, $\beta_2 = [\ce{Ag(NH3)2+}]/([\ce{Ag+}][\ce{NH3}]^2)$.
\textbf{2.} $\log K_1 = 3.32$, $\log K_2 = 7.22 - 3.32 = 3.90$.
\textbf{3.} $[\ce{AgNH3}]^{+}$ predominaria entre pNH3 3.90 e 3.32, um
intervalo que corre ao contrário: ele não tem domínio.
\textbf{4.} \ce{2[AgNH3]+ <=> Ag+ + [Ag(NH3)2]+}, $K = K_2/K_1 = 10^{0.58} =
3.8$.
\textbf{5.} $[\ce{Ag(NH3)2}]^{+}$ abaixo e \ce{Ag+} acima de
$\frac12\log\beta_2 = 3.61$.
\textbf{6.} $\beta_2[\ce{NH3}]^2 = 10^{7.22 - 2.00} = 10^{5.22} = \num{1.7e5}$.
\textbf{7.} $\mathrm{p}K_d = \log\beta_2 = 7.22$.
\textbf{8.} \ce{AgCl(s) + 2NH3 <=> [Ag(NH3)2]+ + Cl-}, $K = \beta_2K_s =
10^{7.22 - 9.75} = 10^{-2.53}$.
\textbf{9.} $s = \sqrt{K}\,[\ce{NH3}] = 10^{-1.265} \times 1.0 =
\qty{0.054}{mol/L}$.
\textbf{10.} Em água, $s = 10^{-4.875} = \qty{1.3e-5}{mol/L}$: cerca de 4000
vezes mais.
\textbf{11.} $0.0543 \times 143.4 = \qty{7.8}{g}$.
\textbf{12.} $[\ce{Ag+}] = K_s/[\ce{Cl-}] = 10^{-9.75}/0.054 =
\qty{3.3e-9}{mol/L} \ll s$.
\textbf{13.} $s = \sqrt{K}(1.0 - 2s)$, $s = 0.0543/(1 + 2 \times 0.0543) =
\qty{0.049}{mol/L}$.
\textbf{14.} A amônia a \qty{1.0}{mol/L} tem pH 11.6, acima de $9.25 + 1$:
só \qty{0.4}{\percent} dela está como \ce{NH4+}.
\textbf{15.} $K = 10^{7.22 - 12.27} = 10^{-5.05}$.
\textbf{16.} $s = 10^{-2.525} = \qty{3.0e-3}{mol/L}$, \qty{0.56}{g/L}.
\textbf{17.} $K = 10^{-8.85}$, $s = 10^{-4.425} = \qty{3.8e-5}{mol/L}$,
\qty{8.8}{mg/L}.
\textbf{18.} Trate cada precipitado com amônia diluída: só o cloreto se
dissolve. A amônia concentrada dissolve o brometo, mais lentamente e em
parte; o iodeto não se dissolve.
\textbf{19.} $s = 1.0/187.8 = \qty{5.3e-3}{mol/L}$; $[\ce{NH3}] = s/\sqrt{K}
= \num{5.3e-3}/10^{-2.525} = \qty{1.8}{mol/L}$.
\textbf{20.} $s = \qty{4.3e-3}{mol/L}$ e $[\ce{NH3}] = \num{4.3e-3}/10^{-4.425}
= \qty{113}{mol/L}$, o dobro da concentração da água na água pura
(\qty{55}{mol/L}): impossível. A amônia não consegue dissolver o iodeto
de prata.
\textbf{21.} \ce{[Ag(NH3)2]+ + Cl- + 2H3O+ -> AgCl(s) + 2NH4+ + 2H2O},
$K = 10^{2.53} \times (10^{9.25})^2 = 10^{21.03}$: o precipitado branco
reaparece, o que confirma o cloreto de prata.
\textbf{22.} $s = 1.0/143.4 = \qty{6.97e-3}{mol/L}$; $[\ce{NH3}] =
s/\sqrt{K} = \num{6.97e-3}/10^{-1.265} = \qty{0.128}{mol/L}$.
\textbf{23.} $2s = \qty{0.0139}{mol/L}$.
\textbf{24.} $[\ce{Ag+}] = 10^{-9.75}/\num{6.97e-3} = \qty{2.6e-8}{mol/L}$,
desprezível diante de $s$.
\textbf{25.} $0.128 + 0.014 = $ \textbf{\qty{0.142}{mol/L} de amônia}:
\qty{0.142}{mol} no litro dissolve exatamente \qty{1.0}{g} de cloreto de
prata.
\end{solution}
